\subsection{What remains of the gap}\label{sec:gap}

How much of the straighter progress of groups is left once the route is the same?

Over a quarter of an hour, group time still progresses straighter than the same route flown alone:
$\Delta\Hloc(1000~\mathrm{s})$ is \IsoDHloc{} \IsoDHlocBand{}, against \IsoBaseDHloc{} for the same pieces
against the solo flights of their site-day and hour (Table~\ref{tab:isolated}). Fixing the route lowers
the gap by \IsoExplained\%, Eq.~\eqref{eq:explained}; its band stays away from zero. The long series,
on the pieces that last long enough, gives about the same gap up to 3000~s, with bands wide enough to
reach zero (Figure~\ref{fig:dhloc}).

\begin{figure}
\centering
\includegraphics[width=0.62\textwidth]{isolated_dhloc.pdf}
\caption{$\Delta\Hloc(\tau)$, group time minus the same route flown alone, for beginners (\IsoPiecesBeg{} pieces),
experienced pilots (\IsoPiecesExp{}) and both (\IsoPieces{}). Continuous: the lags up to 1000~s; dashed: the long series, on the pieces that last
long enough. Bands 5--95\%, thin lines the range when one year is left out.}
\label{fig:dhloc}
\end{figure}

The gap changes sign with the lag: at short lags, within one climb or glide, group time is less straight
than the same route flown alone, and it is more straight from about 200~s on (Figure~\ref{fig:dhloc}).

\begin{table}
\centering\small
\begin{tabular}{@{}l r r c c r@{}}
\toprule
wing level & pieces & matches & against site-day and hour & against isolated, same route & explained \\
\midrule
all & 15\,539 & 85\,524 & $+$0.058 & $+$0.017 [$+$0.011, $+$0.025] & 70\% \\
experienced & 10\,824 & 64\,242 & $+$0.051 & $+$0.016 [$+$0.010, $+$0.024] & 69\% \\
beginners & 4\,715 & 21\,282 & $+$0.068 & $+$0.022 [$+$0.006, $+$0.038] & 68\% \\
\bottomrule
\end{tabular}

\caption{$\Delta\Hloc(1000~\mathrm{s})$ of group time against the solo flights of the same site-day
and hour, and against the isolated solo flights on the same route, with the 5--95\% band; the share
of the first that the route explains, Eq.~\eqref{eq:explained}.}
\label{tab:isolated}
\end{table}

\paragraph{What we keep.} The route, with the season, the hour and the wing level, explains about
\IsoExplained\% of the long-lag gap between group and solo flight. The rest, \IsoDHloc{} at
1000~s, is small and goes with flying together.
